\chapter*{Executive summary}
\addcontentsline{toc}{chapter}{Executive summary}

\paragraph{The question.} In a multi-core quantum computer, moving a logical qubit between cores consumes an entangled (EPR) pair, and every core has only a few
communication memories. When many moves need the same core at once, requests wait: \emph{congestion}. This report documents what was wrong with how the ACE simulator measured that,
how it was corrected, and what the corrected simulator says about ten benchmark circuits on a $4\times4$ mesh of cores.

\paragraph{What was wrong.} ACE split every layer of requests into batches with at most one request per core, so a core never held more than one request and the memory budget could not bind;
the earlier budget study therefore could not measure congestion. That serialization also hid a second fault: two requests that looked alike (same endpoints, window, size and fidelity) were treated as one booking
by SeQUeNCe's reservation code (four requests on one link took up to 68\,ms instead of 2.4\,ms). Ten findings are in Chapter~\ref{ch:audit}.

\paragraph{What was built.} A reservation-identity fix; a per-qubit request manager with event-driven admission (hops of one qubit in order, different qubits concurrent, a request starts only when both of
its cores have a free memory); an event log and an invariant checker; one latency definition valid for every scheduler; and drivers that record every parameter and every run time.
The legacy simulator is untouched and still reproduces its 164-row baseline exactly after every change.

\paragraph{Headline results (Experiments 1 and 2).}
\begin{itemize}
  \item \textbf{The memory budget now binds.} On-demand latency used to be identical at every budget; it now falls with memory. QFT: $14.42$\,ms with one memory per core, $4.72$ with two, $1.59$ with four (the modelled value) and $1.13$ with 32,
        where it saturates (64 memories reproduces 32). The floor of about $1.1$--$1.2$\,ms is generation time.
  \item \textbf{Who pays.} At the modelled four memories the on-demand penalty against saturation is $+40\%$ for QFT, $+12\%$ for QAOA, $+5\%$ for graph state and about $1\%$ or less for the rest. GHZ, Bernstein--Vazirani and BMW never wait.
  \item \textbf{Pre-generation is not free.} Dynamic compiler pre-generation reaches $0.25$\,ms at 32 memories, but its compiler bank takes $\text{budget}-1$ memories, so for the congested workloads it is slower than on-demand until 8 to 16 memories, and it
        lowers the delivered fidelity (QFT down to $0.9185$ and the random circuit, the lowest, to $0.907$, against $0.9498$ on-demand).
  \item \textbf{Where congestion comes from (Experiment 2).} At four memories nearly all overload is \emph{relay} overload: a core is overloaded because routes pass through it, not because routes start or end there. Endpoint demand alone peaks at 4--5, total demand at 24.
        With one memory, own demand matters (QFT 77\% own-overload).
  \item \textbf{Structure predicts cost.} The heatmap's overload share and the simulated on-demand penalty are strongly rank-correlated (Spearman $\rho=0.85$ over 60 points).
\end{itemize}

\paragraph{Is compiler-driven pre-generation good? (Chapters~\ref{ch:why-dynamic} and~\ref{ch:four-policies}.)} \textbf{Against the autonomous protocols CGP and ACGP, yes.} On nine benchmarks, on the mesh and on the torus, dynamic pre-generation has lower latency than both
  at every modelled memory count (the better of CGP and ACGP takes $1.5$ to $2.5$ times as long at four memories), wastes no pair against $34$ to $89\%$ of pairs for CGP and ACGP on the mesh at four memories, and delivers higher fidelity ($0.916$ to $0.927$ against $0.865$ to $0.913$ at four memories).
  \textbf{Against on-demand generation, only partly.} Dynamic is faster on five of the nine benchmarks at four memories and on all of them from sixteen memories on, but on QFT, QAOA, Draper and graph state it is slower at two to four memories (QFT at four: $5.57$ against $1.59$\,ms), and its fidelity is below on-demand's $0.950$ everywhere.
  The earlier simulator showed dynamic ahead because of its serialization (QFT: $41\%$ faster in its own latency, $11\%$ counting the wait for earlier batches). The cause on the corrected simulator was read in the code and measured: on QFT at four memories prepared pairs hold memory cards for $24$\,ms on average
  against $6$\,ms for an ordinary request, a request that has a ready pair still needs a card of its own, and the pair of a later hop of a qubit's chain is made at the start of the previous layer but used only after the hop before it completed. Whether a better plan could avoid this is not established.
  The random circuit was not run for CGP and ACGP.

\paragraph{Beyond Experiments 1 and 2 (shorter chapters).} The torus needs 11--37\% fewer hops, has lower structural overload for five workloads (unchanged for four, higher for graph state); offline re-routing with the mapping fixed removes most of the excess demand and makes on-demand requests up to 42\% faster; ACE already contains a memory cut-off
(\code{cutoff_ratio}) that is unused by default.

\paragraph{What is not settled.} The physical parameters. ACE's coherence time (1\,s) and memory efficiency (0.95) differ from the group's multi-core paper (1\,ms, 0.85). With the paper's values every pre-generated pair expires unused, so every pre-generation
and fidelity result depends on that choice; the congestion results barely do. This needs a decision from the supervisor (Appendix~\ref{app:params}).

\paragraph{How to trust it.} Every result carries controls that must pass (a zero-congestion workload must not respond to the budget; 64 memories must equal 32; completion must be 100\%); the first run failed two of them for non-physical reasons, which were diagnosed
and fixed (Section~\ref{sec:corrections}); every experiment folder records all parameters, the code commit, the machine and the run times. Section~\ref{sec:corrections} lists every mistake made on the way and how it was found.

\paragraph{Status of the evidence.} The historical chapters preserve the earlier
corrected-scheduler and four-policy results, but they are not silently mixed with
the new Golden-study corpus. The Golden study is the current reproducibility
baseline: its complete contract, trace hashes, source lock and acceptance rules
are in Chapter~14. At this report build, the GHZ and BV mesh--XY slices are fully
validated; the remaining base workloads and the separately registered MQT Bench
QPE extension are still executing. The report therefore makes workload-specific
observations, not a claim that compiler-driven Dynamic is universally superior.
